\begin{summarybox}

    \textbf{\textcolor{CSummary}{一条由初中走向高中的理解路线}}

    \textbf{普通三角形：}竖直中线可作两个小三角形的共同底边，
    用两侧的\textbf{水平距离}作高。例：$PM=4$、水平高各为 $2$，
    面积 $=\tfrac12\times4\times2+\tfrac12\times4\times2=8$。

    \textbf{圆的两切线：}两切线段相等，所以中线 $PM$\textbf{垂直}底边 $AB$，
    面积可直接用 $\tfrac12 AB\cdot PM$。

    \textbf{抛物线的两切线：}中线 $PM$\textbf{平行对称轴}，通常\textbf{不垂直} $AB$；
    因此回到第一步的分割方法，得到
    $S=\tfrac12 PM\cdot|x_1-x_2|$。
    \textbf{先分清中线与高，再选择正确的“底和高”——这是本结论最值得迁移的方法。}

    \medskip

    \textbf{\textcolor{CSummary}{一句话核心}}

    \textbf{两切线交点 $P$ 与切点弦中点 $M$ 处于同一条竖直线上。}
    所以不用求斜底边的长度，只要算出竖直中线 $PM$，
    就能用初中的三角形面积公式得到三次方的面积结论。

    \medskip
    \textbf{\textcolor{CSummary}{使用条件}}

    抛物线 $x^2=2py\;(p>0)$；$A,B$ 为曲线上两个\textbf{不同}的点；
    过 $A,B$ 的切线交于 $P$，$M$ 是 $AB$ 的中点。

    \medskip
    \textbf{\textcolor{CSummary}{四个必记等式}}
    \[
    \boxed{\begin{gathered}
        P\left(\frac{x_1+x_2}{2},\frac{x_1x_2}{2p}\right),\\[3pt]
        PM=\frac{(x_1-x_2)^2}{4p},\qquad PM\parallel y\text{轴},\\[3pt]
        S_{\triangle PAB}=\frac12 PM\cdot|x_1-x_2|
        =\frac{|x_1-x_2|^3}{8p},\\[3pt]
        P=(x_0,y_0)\ \Longrightarrow\ AB:\ x_0x=p(y+y_0).
    \end{gathered}}
    \]

    \medskip
    \textbf{\textcolor{CSummary}{四种证明的使用次序}}
    \begin{itemize}[leftmargin=*, itemsep=0.35em, topsep=0.3em]
        \item \textbf{先理解：}切线点式 $\to$ 横坐标取平均 $\to$ 竖直中线 $\to$ 两块三角形面积。
        \item \textbf{考场最快：}相减求 $x_P$，回代求 $y_P$，用平方差算 $PM$，直接出面积。
        \item \textbf{步骤最少：}设 $x_1=x_0-t,x_2=x_0+t$，两大步求顶点、中点与面积。
        \item \textbf{方法迁移：}通过切线方程、弦长、点到直线距离求面积。
    \end{itemize}

    \medskip
    \par\noindent\begin{minipage}{\linewidth}
    \textbf{\textcolor{CSummary}{最有价值的几何联系}}

    若 $N$ 是竖直线 $PM$ 与抛物线的交点，则
    \[
        \boxed{PN=NM=\frac{(x_1-x_2)^2}{8p}}.
    \]
    也就是 $N$ 为\textbf{线段 $PM$ 的中点}。
    “水平宽度的一次方 $\times$ 竖直中线的平方关系”解释了面积公式中的三次方。
    \end{minipage}\par

    \medskip
    \textbf{\textcolor{CSummary}{三个考前检查}}
    \begin{itemize}[leftmargin=*, itemsep=0.3em, topsep=0.3em]
        \item 绝对值不能丢：面积 $S=|x_1-x_2|^3/(8p)>0$。
        \item $PM$ 是中线，\textbf{通常不是} $AB$ 上的高。
        \item 已知交点 $P(x_0,y_0)$ 时，先确认 $x_0^2-2py_0>0$ 才能作出两个不同实切点。
    \end{itemize}
\end{summarybox}
